\[ S = (e1=60, f1=61, g1=62, h1=63),\quad T_S = \begin{pmatrix}0 & 0 & 0 & 0 \\ 0 & 0 & 0 & 1 \\ 1 & 0 & 0 & 0 \\ 0 & 0 & 0 & 0\end{pmatrix},\quad b_S = \begin{pmatrix}6 \\ 0 \\ 0 \\ 4\end{pmatrix}\;\longmapsto\; (Tb)_S =\begin{pmatrix}0 \\ 4 \\ 6 \\ 0\end{pmatrix} \]
\[ R_S = \begin{pmatrix}0 & 0 & 1 & 0 \\ 0 & 0 & 0 & 0 \\ 0 & 0 & 0 & 0 \\ 0 & 1 & 0 & 0\end{pmatrix},\qquad (R\,b')_S = \begin{pmatrix}6 \\ 0 \\ 0 \\ 4\end{pmatrix} \]
\begin{itemize}
\item \textbf{Coups de la pièce.} Le masque $M$ contient 2 destinations, et la pièce a 3 coups légaux dans cette position.
\item \textbf{Matrice.} Ici, $R$ rend $b_S$ sans perte. La matrice $T$ est de rang 62 sur 64.
\item \textbf{Forme réduite.} La ligne $(f, t, k, c, f', t') = (60, 62, 1, 64, 63, 61)$ donne les 4 écritures $b'_{62} = b_{60},\ b'_{60} = 0,\ b'_{61} = b_{63},\ b'_{63} = 0$.
\item \textbf{Plan.} Le coup appartient au plan $p = 68$ (« O-O »), de décalage $\delta_p = +2$, et $t = f + \delta_p = 60 +2 = 62$.
\end{itemize}
